\chapter*{TÓM LƯỢC}
\addcontentsline{toc}{chapter}{TÓM LƯỢC}

Giữa nhịp sống hối hả của đô thị hiện đại, khi con người ngày càng khao khát tìm về những phút giây an nhiên và lối sống chánh niệm, văn hóa thưởng trà truyền thống đang dần chuyển mình để tiếp cận thế hệ người tiêu dùng mới. Nắm bắt xu hướng giao thoa giữa nét đẹp mộc mạc của văn hóa Việt và sức mạnh của công nghệ số, đồ án \textbf{"Hệ thống thương mại điện tử kinh doanh trà và phụ kiện pha trà"} được thực hiện nhằm kiến tạo một không gian trải nghiệm và mua sắm trực tuyến tinh tế mang thương hiệu \textbf{Sen Việt Tea}.

Đồ án là bức tranh toàn cảnh về quy trình chuyển đổi số cho một thương hiệu ngành F\&B đặc thù, từ bước phác thảo chân dung người tiêu dùng trẻ bận rộn, thiết kế hệ thống nhận diện (Design System) tinh giản, đến khâu lập trình và triển khai kỹ thuật thực tế. Bằng việc tận dụng linh hoạt lõi thương mại điện tử WooCommerce trên nền tảng mã nguồn mở WordPress, kết hợp cùng giao diện Flatsome và các giải pháp tối ưu hóa hiệu năng, hệ thống không chỉ giải quyết triệt để bài toán bán hàng tự động mà còn đáp ứng những tiêu chuẩn khắt khe về bảo mật và thân thiện với công cụ tìm kiếm (SEO).

Kết quả mang lại là một website thương mại điện tử hoàn thiện, nơi khách hàng không chỉ dễ dàng thực hiện các thao tác từ xem sản phẩm, quản lý giỏ hàng đến thanh toán trực tuyến mượt mà, mà còn được đắm chìm vào câu chuyện văn hóa nghệ thuật trà đạo qua từng khung hình, góc chữ. Đồ án minh chứng cho việc ứng dụng công nghệ phần mềm không hề làm khô cứng đi những giá trị truyền thống, mà trái lại, còn là bệ phóng vững chắc giúp các doanh nghiệp kinh doanh sản phẩm văn hóa sẵn sàng vươn xa trên môi trường số.
